% --------------------------------
% Estilo dos títulos do CAPÍTULOS
% --------------------------------
\makechapterstyle{mymadsen}{%
  
  \chapterstyle{default}
  
  \renewcommand*{\chapnamefont}
  {%
    \normalfont\large\scshape\hspace{12cm}
    }
  
  \renewcommand*{\chaptitlefont}
  {%
    \normalfont\Huge\bfseries%\sffamily
    \raggedleft
    }
  
  \renewcommand*{\chapternamenum}{}
  
  \setlength{\chapindent}{0.08 in} %antes era {0.75 in}
  
  \setlength\beforechapskip{0.01 in}
  
  \renewcommand*{\printchapternum}
  {%
    \makebox[0pt][l]{\hspace{0.5cm}%
      \resizebox{!}{1cm}{%
        \bfseries%\sffamily
        \thechapter}%
    }%
  }%
  
  \renewcommand*{\printchapternonum}{%
    \chapnamefont \phantom{\printchaptername \chapternamenum%
      \makebox[0pt][l]{\hspace{0.4em}%
        \resizebox{!}{4ex}{%
          \chapnamefont\bfseries%\sffamily 
          10}%
      }%
    }%
  }%
  
  \setlength{\beforechapskip}{-1.5cm}
  
  \setlength\midchapskip{0cm}
  
  \renewcommand*{\printchaptertitle}[1]{%
    \begin{adjustwidth}{-\chapindent}{}
      \par\hrulefill\vskip\midchapskip
      \raggedleft \chaptitlefont ##1\par\nobreak
    \end{adjustwidth}}
}




% --------------------------------------
%   Estabelecendo o estilo do capítulo
% --------------------------------------
    \chapterstyle{mymadsen}
% --------------------------------------




%---------------------------------
%	Identação de Parágrafos
%---------------------------------

% O tamanho da identação do parágrafo é dado por:
\setlength{\parindent}{1.3cm}

% Controle do espaçamento entre um parágrafo e outro:
\setlength{\parskip}{0.2cm}  % tente também \onelineskip